\chapter{Electric, Magnetic, and Electromagnetic Fields}
\label{ch:fields}

Everything a radio does --- store energy in a capacitor, store energy in
an inductor, radiate a signal from an antenna, propagate that signal
across an ocean --- is a manifestation of electric and magnetic fields.
This chapter builds the field concepts needed to understand components
(Part~II) and antennas and propagation (Part~IV).

\section{The Electric Field}

An \textbf{electric field} exists in the region around any electric
charge; it is defined as the force per unit charge that a small positive
``test charge'' would experience at a given point:
\[
E = \frac{F}{q} \qquad [\text{volts per metre, V/m}]
\]
Field lines point away from positive charges and toward negative charges.
Between two flat, oppositely charged parallel plates separated by
distance $d$ with potential difference $U$ between them, the field is
uniform and given by
\[
E = \frac{U}{d}.
\]
This is precisely the geometry inside a parallel-plate capacitor
(Chapter~\ref{ch:capacitors}): a capacitor is, at its core, a device for
storing energy in an electric field confined between two conductors.

\begin{keyidea}
The electric field describes how strongly a charge distribution can push
or pull on other charges. It is the mechanism by which a capacitor stores
energy and by which a transmitting antenna's near field couples to
nearby objects.
\end{keyidea}

\section{The Magnetic Field}

A \textbf{magnetic field} is produced by moving charge (current) and is
described by the flux density $B$, measured in \textbf{tesla (T)}. Unlike
electric field lines, magnetic field lines have no starting or ending
point; they form closed loops. For a long, straight conductor carrying
current $I$, the field at perpendicular distance $r$ is
\[
B = \frac{\mu_0 I}{2\pi r}
\]
where $\mu_0 = 4\pi \times 10^{-7}\,\text{H/m}$ is the \textbf{permeability
of free space}. The direction of the field circles the conductor
according to the \textbf{right-hand rule}: point the thumb of your right
hand in the direction of conventional current flow, and your fingers curl
in the direction of the magnetic field.

Winding a conductor into a coil concentrates and reinforces this field
inside the loops, which is precisely how an \textbf{inductor}
(Chapter~\ref{ch:inductors}) stores energy in a magnetic field, and how a
loop antenna couples to the magnetic component of a radio wave.

\begin{center}
\begin{tikzpicture}[scale=0.9]
\draw[thick,-{Latex}] (0,0) -- (3,0) node[right] {$I$};
\foreach \r in {0.5,0.9,1.3}{
  \draw[thick,-{Latex}] (1.5,0) circle (\r);
}
\node at (1.5,-1.8) {Field circles the conductor (right-hand rule)};
\end{tikzpicture}
\end{center}

\section{Magnetic Flux and Faraday's Law}

\textbf{Magnetic flux}, $\Phi$, measured in \textbf{webers (Wb)}, is the
total magnetic field passing through a given area:
\[
\Phi = B \cdot A
\]
for a uniform field $B$ perpendicular to area $A$. \textbf{Faraday's law of
induction} states that a changing magnetic flux through a loop induces an
EMF around that loop:
\[
\mathcal{E} = -N\frac{d\Phi}{dt}
\]
where $N$ is the number of turns of the loop and the negative sign (Lenz's
law) indicates that the induced EMF always opposes the change that
created it. Faraday's law is the working principle behind
\textbf{inductors} (an inductor opposes changes in current because that
current change causes a changing flux, which induces an opposing EMF),
\textbf{transformers} (Chapter~\ref{ch:transformers}), and the reception
of radio signals by any loop or wire antenna.

\begin{keyidea}[Lenz's law, intuitively]
An induced current always flows in the direction that opposes the change
that produced it. This is why inductors resist sudden changes in
current, and why energy must be supplied to overcome this opposition when
current is increased or decreased quickly.
\end{keyidea}

\section{The Electromagnetic Field and Maxwell's Insight}

Electric and magnetic fields are not independent phenomena; they are two
aspects of a single \textbf{electromagnetic field}. James Clerk Maxwell
showed in the 1860s, through four coupled equations now called
\textbf{Maxwell's equations}, that:

\begin{itemize}
\item A changing electric field produces a magnetic field (this
completes the symmetry of Faraday's law, and was Maxwell's key
addition --- the ``displacement current'' term).
\item A changing magnetic field produces an electric field (Faraday's
law).
\end{itemize}

These two facts together mean that, once created, a time-varying
electric field and a time-varying magnetic field can regenerate each
other indefinitely, propagating outward from their source as a
self-sustaining wave --- an \textbf{electromagnetic wave} --- without
needing a material medium to carry them. Maxwell calculated the
predicted speed of such a wave from purely electrical and magnetic
constants,
\[
c = \frac{1}{\sqrt{\mu_0 \varepsilon_0}} \approx 2.998 \times 10^{8}\,\text{m/s},
\]
and found it matched the already-measured speed of light --- correctly
concluding that light itself is an electromagnetic wave. Radio waves,
microwaves, infrared, visible light, ultraviolet, X-rays, and gamma rays
are all electromagnetic waves, differing only in frequency (and hence
wavelength); amateur radio operates in the radio-frequency portion of
this same spectrum.

\section{Properties of a Traveling EM Wave}

In a radio wave traveling through free space, the electric field $E$ and
magnetic field $H$ oscillate perpendicular to each other and to the
direction of propagation, in phase, with a fixed ratio between them (the
\textbf{impedance of free space}, $Z_0 = \sqrt{\mu_0/\varepsilon_0}
\approx 377\,\Omega$).

\begin{center}
\begin{tikzpicture}[scale=0.85]
\draw[thick,blue,domain=0:6,samples=60] plot ({\x},{1.1*sin(\x*90)});
\node[blue,anchor=south] at (1.5,1.55) {$E$ (vertical plane)};
\foreach \x in {0.5,1.5,2.5,3.5,4.5,5.5}{
  \draw[thick,red] (\x,0) ellipse (0.06 and 0.35);
}
\node[red,anchor=north] at (4.5,-0.5) {$H$ (horizontal plane)};
\draw[-{Latex},thick] (-0.3,-1.7) -- (6.3,-1.7);
\node[anchor=north] at (3,-1.75) {direction of travel};
\end{tikzpicture}
\end{center}

The \textbf{wavelength} $\lambda$, \textbf{frequency} $f$, and
\textbf{speed of light} $c$ are related by
\[
\boxed{c = f\lambda} \qquad \Longrightarrow \qquad \lambda = \frac{c}{f}, \quad f = \frac{c}{\lambda}.
\]
This single relationship is one of the most exam-relevant formulas in the
entire book: it converts between the two ways radio frequencies are
described (as a frequency in Hz, or as a wavelength in metres, as in
``the 20-metre band'') and underlies antenna dimension calculations in
Chapter~\ref{ch:antennas}.

\begin{examplebox}[Frequency--wavelength conversion]
Find the wavelength corresponding to 14.200\,MHz.

\emph{Solution.}
\[
\lambda = \frac{c}{f} = \frac{2.998\times10^{8}\,\text{m/s}}{14.200\times10^{6}\,\text{Hz}} \approx 21.1\,\text{m}.
\]
This is why 14\,MHz is colloquially called ``the 20\,m band'' --- 21.1\,m
rounds, by amateur radio band-naming convention, to the nearest
traditional metre-band label.
\end{examplebox}

\section{Polarization}

The \textbf{polarization} of an EM wave describes the orientation of its
electric field vector. If $E$ stays in a fixed plane as the wave travels
(e.g., always vertical, or always horizontal), the wave is
\textbf{linearly polarized}; a horizontal dipole radiates a horizontally
polarized wave, a vertical antenna radiates a vertically polarized wave.
If the $E$ field instead rotates as the wave travels, tracing a helix, the
wave is \textbf{circularly} (or, more generally, \textbf{elliptically})
polarized --- used heavily in satellite work, where a spacecraft's
orientation relative to a ground station is constantly changing, so
neither pure horizontal nor pure vertical polarization stays aligned.
Polarization mismatch between a transmitting and receiving antenna causes
signal loss (Chapter~\ref{ch:antennas}), which is why matching
polarization matters when choosing or orienting an antenna.

\begin{quickreview}
\item Electric field $E = F/q$ (V/m); magnetic flux density $B$ (tesla);
flux $\Phi = BA$ (weber).
\item Faraday's law: a changing flux induces an EMF, $\mathcal{E} =
-N\,d\Phi/dt$; Lenz's law fixes the sign (opposition to change).
\item Maxwell's equations show that changing $E$ and $B$ fields sustain
each other, propagating as an electromagnetic wave at $c \approx
3\times10^{8}\,\text{m/s}$.
\item $c = f\lambda$ links frequency and wavelength --- essential for
antenna design and band naming.
\item Polarization describes the orientation of the wave's electric
field; mismatched polarization causes signal loss.
\end{quickreview}

\begin{practiceproblems}
\item Find the magnetic field 5\,cm from a straight wire carrying 2\,A.
\item A single-turn loop of area 0.01\,m$^2$ experiences a flux change
from 0 to 5\,mWb in 0.1\,s. Find the induced EMF.
\item Calculate the wavelength for 145.500\,MHz and for 433.500\,MHz;
identify the traditional amateur band name for each.
\item Explain, in terms of Maxwell's equations, why a radio wave can
cross the vacuum of space with no medium to carry it, unlike a sound
wave.
\item A satellite station uses circular polarization while a fixed home
station uses a horizontal linear-polarized antenna. Explain qualitatively
why using circular polarization at both ends often improves the link.
\end{practiceproblems}
